\documentclass[11pt]{article}

\usepackage[a4paper,top=22mm,bottom=25mm,left=22mm,right=22mm]{geometry}
\usepackage{amsmath,amssymb,amsthm,mathtools}
\usepackage{enumitem}
\usepackage[colorlinks=true,linkcolor=blue,citecolor=blue,urlcolor=blue]{hyperref}

\mathtoolsset{showonlyrefs=true}
\setlength{\emergencystretch}{2em}

\newtheorem{theorem}{Theorem}[section]
\newtheorem{lemma}[theorem]{Lemma}
\newtheorem{proposition}[theorem]{Proposition}
\newtheorem{corollary}[theorem]{Corollary}
\newtheorem{assumption}[theorem]{Assumption}
\newtheorem{conjecture}[theorem]{Conjecture}

\newcommand{\E}{\mathbb{E}}
\newcommand{\Var}{\operatorname{Var}}
\newcommand{\Cov}{\operatorname{Cov}}
\newcommand{\Tr}{\operatorname{Tr}}
\newcommand{\tr}{\operatorname{tr}}
\newcommand{\N}{\mathcal{N}}
\newcommand{\I}{\mathbb{I}}
\newcommand{\C}{\mathbb{C}}
\newcommand{\R}{\mathbb{R}}
\newcommand{\calI}{\mathcal{I}}
\newcommand{\calG}{\mathcal{G}}
\newcommand{\dd}{\,d}
\newcommand{\ind}{\mathbf{1}}
\newcommand{\norm}[1]{\left\lVert #1\right\rVert}
\newcommand{\abs}[1]{\left\lvert #1\right\rvert}
\newcommand{\HS}{\mathrm{HS}}
\newcommand{\Comp}{\operatorname{Comp}}

\title{A high-temperature CLT for the SYK free energy: a two-replica cluster-expansion reduction}
\author{}
\date{}

\begin{document}

\maketitle
\paragraph{Status.}
The result of Gamarnik--Pernice--Schmidhuber--Zlokapa gives the high-temperature
free-energy limit by a sparse-hypergraph and saddle-point analysis, while
Anschuetz--Chen--Kiani--King give an annealed/quenched comparison strong enough
for the LLN level.  The proof below uses these only as the locally uniform LLN
input.  The CLT-level reduction is the Gaussian-chaos analysis.  The remaining
input is reduced to a two-point, rather than all-orders, marked estimate: a
uniform $L^2$ bound on the Kubo covariance matrix of two SYK monomials.  This is
stated in Assumption~\ref{ass:two-marked-susceptibility}.


\tableofcontents

\section{Introduction and setup}
\label{sec:introduction-setup}

\subsection{Reader's guide and dependency of assumptions}

This note is meant to separate the purely Gaussian-chaos part of the CLT from
one genuinely high-temperature cluster-expansion input.  The only input which is
not proved in the present note is the resummed two-replica $L^2$ cluster
expansion, Assumption~\ref{ass:two-replica-cluster-expansion}.  Once this
assumption is supplied, the rest of the CLT proof is self-contained.

The logical dependencies are as follows.  The LLN input is
Assumption~\ref{ass:local-lln}.  The remaining CLT-level input is first stated in
the softer form of the two-marked susceptibility estimate,
Assumption~\ref{ass:two-marked-susceptibility}.  Section~\ref{sec:two-marked-proof-attempt}
then reduces this susceptibility estimate to the more structural cluster
assumption, Assumption~\ref{ass:two-replica-cluster-expansion}.  In symbols,
\begin{equation}\label{eq:reader-main-dependency}
  \begin{gathered}
    \text{Assumption~\ref{ass:local-lln}}
    +
    \text{Assumption~\ref{ass:two-replica-cluster-expansion}}
    \\
    \Longrightarrow
    \text{Assumption~\ref{ass:two-marked-susceptibility}}
    \Longrightarrow
    \text{Theorem~\ref{thm:main-clt}}.
  \end{gathered}
\end{equation}
Equivalently, after the LLN input, the only genuinely unproved input is
Assumption~\ref{ass:two-replica-cluster-expansion}.

More explicitly, the proof is organized through the following chain.
\begin{equation}\label{eq:reader-chain}
  \begin{gathered}
    \text{resummed two-replica }L^2\text{ cluster expansion}
    \\
    \Longrightarrow
    \text{two-marked susceptibility}
    \Longrightarrow
    \text{Hilbert--Schmidt Hessian bound}
    \\
    \Longrightarrow
    \text{infinite chaos tail closure}
    \Longrightarrow
    \text{rank-one second-chaos CLT}.
  \end{gathered}
\end{equation}

The parts proved in this note are the following.
\begin{enumerate}[label=\textup{(\roman*)},leftmargin=*]
  \item The derivative convergence
  $N^{-1}\partial_\beta\E L_N(\beta)\to\ell_q'(\beta)$ is derived from the
  locally uniform LLN by convexity; see Proposition~\ref{prop:derivative-lln-from-lln}.

  \item The first chaos vanishes by Majorana sign-gauge symmetry, and the second
  chaos is computed exactly; see Lemma~\ref{lem:selection},
  Lemma~\ref{lem:second-chaos}, and Proposition~\ref{prop:exact-reduction}.

  \item The second chaos gives the whole limiting Gaussian field once the higher
  chaoses are negligible; see Lemma~\ref{lem:p2-clt}.

  \item The two-marked susceptibility implies a Hilbert--Schmidt bound on the
  Hessian, and this yields the infinite chaos tail closure by the Malliavin
  identity; see Lemma~\ref{lem:kubo-hessian},
  Lemma~\ref{lem:malliavin-second-tail}, and
  Proposition~\ref{prop:higher-from-two-marked}.

  \item The fixed finite chaoses are negligible by a finite marked-diagram
  counting estimate; see Proposition~\ref{prop:finite-chaos-closure}.

  \item The zero-source Duhamel cumulants admit a dependency-graph tree bound,
  and the connected parity counting lemma is proved; see
  Proposition~\ref{prop:duhamel-tree-graph} and
  Lemma~\ref{lem:connected-parity-counting}.
\end{enumerate}

The parts not proved in this note are deliberately isolated.  The naive
single-replica activity expansion fails because the absolute activity has size
$N^{(q-1)/2}$ near a marked edge.  The naive coefficientwise tree expansion also
fails because repeated hyperedges create occurrence-level entropy which should
be resummed.  These failures are recorded in
Subsection~\ref{subsec:why-naive-cluster-fails}.  Thus the correct missing input
is not a coefficientwise tree estimate, but the resummed two-replica estimate of
Assumption~\ref{ass:two-replica-cluster-expansion}.

For a first reading, one can read the setup and
Section~\ref{sec:clt-reduction}, then read the statement of
Assumption~\ref{ass:two-replica-cluster-expansion} and
Proposition~\ref{prop:cluster-expansion-to-susceptibility}.  The remaining part
of Section~\ref{sec:two-marked-proof-attempt} explains why this assumption is
the natural high-temperature input and why the more naive alternatives are too
strong.

\subsection{Model and main theorem}

Fix an even integer $q\geq4$.  Let $N$ be even and let
$\psi_1,\ldots,\psi_N$ be Majorana operators on $\C^{2^{N/2}}$ satisfying
\begin{equation}\label{eq:majorana-relations}
  \psi_i\psi_j+\psi_j\psi_i=2\delta_{ij}\I.
\end{equation}
Thus $\psi_i^2=\I$.  For $I=\{i_1<\cdots<i_q\}\subset[N]$, set
\begin{equation}\label{eq:psiI}
  \psi_I=i^{q/2}\psi_{i_1}\cdots\psi_{i_q}.
\end{equation}
This is the paper notation for the Lean term
$(\operatorname{phase} q)\cdot M.\operatorname{psiProduct}(I)$.  Then
$\psi_I$ is Hermitian and $\psi_I^2=\I$.

Let
\begin{equation}\label{eq:calI-Nn}
  \calI_N=\{I\subset[N]: |I|=q\},
  \qquad
  M_N=|\calI_N|=\binom{N}{q}.
\end{equation}
Thus $\calI_N$ is the paper notation for Lean's
$\operatorname{InteractionIndex}\,N\,q$.
Let $(g_I)_{I\in\calI_N}$ be i.i.d. standard Gaussian random variables and set
\begin{equation}\label{eq:sigma}
  \sigma_N^2=\frac{(q-1)!}{N^{q-1}}.
\end{equation}
The Gaussian SYK Hamiltonian is
\begin{equation}\label{eq:Hn}
  H_N(g)=\sigma_N\sum_{I\in\calI_N}g_I\psi_I.
\end{equation}
Equivalently, this is the Lean Hamiltonian with coupling
$J_I=\sigma_Ng_I$.  We use normalized trace $\tr=2^{-N/2}\Tr$ and define
\begin{equation}\label{eq:Zn-Ln}
  Z_N(\beta)=\tr\exp(-\beta H_N(g)),
  \qquad
  L_N(\beta)=\log Z_N(\beta).
\end{equation}

For an observable $A$, write
\begin{equation}\label{eq:thermal-average}
  \langle A\rangle_{\beta,g}
  =
  \frac{\tr(Ae^{-\beta H_N(g)})}{\tr(e^{-\beta H_N(g)})}.
\end{equation}
For $I,J\in\calI_N$, define the Kubo connected covariance
\begin{equation}\label{eq:kubo-covariance}
  \mathcal C_{I,J}(\beta,g)
  =
  \int_0^1
  \left\{
    \frac{
      \tr\left(\psi_I e^{-s\beta H_N(g)}
      \psi_J e^{-(1-s)\beta H_N(g)}\right)
    }{Z_N(\beta)}
    -
    \langle\psi_I\rangle_{\beta,g}
    \langle\psi_J\rangle_{\beta,g}
  \right\}
  \dd s.
\end{equation}
The rescaled centered fluctuation is
\begin{equation}\label{eq:Xn}
  X_N(\beta)=N^{q/2-1}\bigl(L_N(\beta)-\E L_N(\beta)\bigr).
\end{equation}
Let $P_r$ be the orthogonal projection in $L^2$ of the Gaussian disorder onto
the $r$-th homogeneous Wiener chaos.

The proof uses the following LLN input from the high-temperature free-energy
analysis.  The derivative part is not assumed separately; it is proved in
Proposition~\ref{prop:derivative-lln-from-lln}.

\begin{assumption}[Locally uniform high-temperature LLN]\label{ass:local-lln}
  There exist $\beta_0=\beta_0(q)>0$ and an analytic function
  $\ell_q:(-\beta_0,\beta_0)\to\R$ such that, for every compact
  $K\subset(-\beta_0,\beta_0)$,
  \begin{equation}\label{eq:local-lln}
    \sup_{\beta\in K}
    \left|
      \frac1N\E L_N(\beta)-\ell_q(\beta)
    \right|\to0.
  \end{equation}
\end{assumption}

The CLT-level input is the following high-temperature two-marked susceptibility
estimate.  Section~\ref{sec:two-marked-proof-attempt} reduces this estimate to
a single explicit, plausible, but presently unproved input: the resummed
two-replica $L^2$ cluster-expansion estimate in
Assumption~\ref{ass:two-replica-cluster-expansion}.

\begin{assumption}[Two-marked high-temperature susceptibility]
  \label{ass:two-marked-susceptibility}
  There exists $\beta_{\rm sus}=\beta_{\rm sus}(q)>0$ such that, for every
  compact $K\subset(-\beta_{\rm sus},\beta_{\rm sus})$, the quantity
  \begin{equation}\label{eq:two-marked-susceptibility}
    S_{q,K}
    :=
    \sup_{N\geq1}
    \sup_{\beta\in K}
    \sup_{I\in\calI_N}
    \sum_{J\in\calI_N}
    \E\left[
      \left|
        \mathcal C_{I,J}(\beta,g)
      \right|^2
    \right]
    <\infty.
  \end{equation}
\end{assumption}

\begin{theorem}[High-temperature CLT from two-marked susceptibility]
  \label{thm:main-clt}
  Assume Assumptions~\ref{ass:local-lln} and
  \ref{ass:two-marked-susceptibility}.  Let
  \begin{equation}\label{eq:beta-star}
    \beta_*<\min\{\beta_0,\beta_{\rm sus}\}.
  \end{equation}
  Then, for every finite set $\beta_1,\ldots,\beta_k\in(-\beta_*,\beta_*)$,
  \begin{equation}\label{eq:fdd-clt}
    (X_N(\beta_1),\ldots,X_N(\beta_k))
    \Rightarrow
    \left(
      \sqrt{\frac{q!}{2}}\,\beta_1\ell_q'(\beta_1)\xi,
      \ldots,
      \sqrt{\frac{q!}{2}}\,\beta_k\ell_q'(\beta_k)\xi
    \right),
  \end{equation}
  where $\xi\sim\N(0,1)$.  Equivalently, the limiting covariance is
  \begin{equation}\label{eq:limit-covariance}
    \Cov(\calG_q(\beta),\calG_q(\beta'))
    =
    \frac{q!}{2}\beta\ell_q'(\beta)\beta'\ell_q'(\beta').
  \end{equation}
\end{theorem}

\section{CLT reduction}
\label{sec:clt-reduction}

\subsection{Derivative LLN from convexity}

\begin{lemma}[Convex convergence implies derivative convergence]
  \label{lem:convex-derivative}
  Let $I\subset\R$ be an open interval.  Let $f_N:I\to\R$ be differentiable
  convex functions and assume that $f_N\to f$ locally uniformly on $I$, where
  $f\in C^1(I)$.  Then, for every compact $K\subset I$,
  \begin{equation}\label{eq:convex-derivative-conclusion}
    \sup_{x\in K}|f_N'(x)-f'(x)|\to0.
  \end{equation}
\end{lemma}

\begin{proof}
  Fix a compact $K\subset I$.  Choose $h_0>0$ such that
  $K+[-h_0,h_0]\subset I$.  For $0<h<h_0$ and $x\in K$, convexity gives
  \begin{equation}\label{eq:convex-difference-squeeze}
    \frac{f_N(x)-f_N(x-h)}{h}
    \leq f_N'(x)
    \leq
    \frac{f_N(x+h)-f_N(x)}{h}.
  \end{equation}
  Put
  \begin{equation}\label{eq:epsilon-nh}
    \epsilon_N(h)=\sup_{y\in K+[-h,h]}|f_N(y)-f(y)|.
  \end{equation}
  Then $\epsilon_N(h)\to0$ for fixed $h$.  From
  \eqref{eq:convex-difference-squeeze},
  \begin{equation}\label{eq:upper-convex-derivative}
    f_N'(x)-f'(x)
    \leq
    \frac{f(x+h)-f(x)}{h}-f'(x)
    +\frac{2\epsilon_N(h)}{h},
  \end{equation}
  and
  \begin{equation}\label{eq:lower-convex-derivative}
    f_N'(x)-f'(x)
    \geq
    \frac{f(x)-f(x-h)}{h}-f'(x)
    -\frac{2\epsilon_N(h)}{h}.
  \end{equation}
  Since $f'$ is uniformly continuous on $K+[-h_0,h_0]$, the deterministic
  difference-quotient errors converge to zero uniformly on $K$ as $h\downarrow0$.
  First choose $h$ and then choose $N$ to make $\epsilon_N(h)/h$ small.  This
  proves the claim.
\end{proof}

\begin{proposition}[Derivative-strength LLN]
  \label{prop:derivative-lln-from-lln}
  Under Assumption~\ref{ass:local-lln}, for every compact
  $K\subset(-\beta_0,\beta_0)$,
  \begin{equation}\label{eq:derivative-lln}
    \sup_{\beta\in K}
    \left|
      \frac1N\partial_\beta\E L_N(\beta)-\ell_q'(\beta)
    \right|
    \to0.
  \end{equation}
\end{proposition}

\begin{proof}
  Let $f_N(\beta)=N^{-1}\E L_N(\beta)$.  For fixed disorder,
  $\beta\mapsto L_N(\beta)$ is convex, because $L_N''(\beta)$ is the thermal
  variance of $H_N(g)$.  Therefore $f_N$ is convex.  Differentiability follows
  by differentiating under the Gaussian expectation; this is justified by
  $\norm{H_N(g)}\leq \sigma_N\sum_{I\in\calI_N}|g_I|$ and Gaussian
  exponential integrability in finite dimension.  Assumption~\ref{ass:local-lln} gives
  locally uniform convergence of $f_N$ to the analytic, hence $C^1$, function
  $\ell_q$.  Lemma~\ref{lem:convex-derivative} gives
  \eqref{eq:derivative-lln}.
\end{proof}

\subsection{Hermite decomposition and gauge selection}

Let $\gamma_N$ be the product Gaussian law of $(g_I)_{I\in\calI_N}$.  For a
multi-index $\alpha=(\alpha_I)_{I\in\calI_N}$ define
\begin{equation}\label{eq:multi-index}
  |\alpha|=\sum_{I\in\calI_N}\alpha_I,
  \qquad
  \alpha!=\prod_{I\in\calI_N}\alpha_I!,
  \qquad
  H_\alpha(g)=\prod_{I\in\calI_N}h_{\alpha_I}(g_I),
\end{equation}
where $h_k$ is the probabilists' Hermite polynomial.

\begin{lemma}[Hermite coefficient formula]\label{lem:hermite-coeff}
  For every fixed $\beta\in\R$,
  \begin{equation}\label{eq:chaos-decomposition}
    L_N(\beta)-\E L_N(\beta)
    =
    \sum_{r\geq1}P_rL_N(\beta)
  \end{equation}
  in $L^2(\gamma_N)$, and
  \begin{equation}\label{eq:hermite-coeff}
    P_rL_N(\beta)
    =
    \sum_{|\alpha|=r}c_{N,\alpha}(\beta)H_\alpha(g),
    \qquad
    c_{N,\alpha}(\beta)
    =
    \frac1{\alpha!}\E[\partial^\alpha L_N(\beta)].
  \end{equation}
\end{lemma}

\begin{proof}
  The map $g\mapsto L_N(\beta)$ is $C^\infty$ and has at most exponential
  growth in the finite-dimensional Gaussian vector $g$.  Hence its Hermite
  expansion converges in $L^2(\gamma_N)$.  Orthogonality gives
  \begin{equation}\label{eq:hermite-projection}
    P_rL_N(\beta)
    =
    \sum_{|\alpha|=r}
    \frac{\E[L_N(\beta)H_\alpha(g)]}{\alpha!}H_\alpha(g).
  \end{equation}
  Repeated Gaussian integration by parts gives
  \begin{equation}\label{eq:gibp-hermite}
    \E[L_N(\beta)H_\alpha(g)]=\E[\partial^\alpha L_N(\beta)],
  \end{equation}
  proving \eqref{eq:hermite-coeff}.
\end{proof}

For $s=(s_1,\ldots,s_N)\in\{\pm1\}^N$ define
\begin{equation}\label{eq:chi-s}
  \chi_s(I)=\prod_{a\in I}s_a,
  \qquad
  (\chi_sg)_I=\chi_s(I)g_I.
\end{equation}

\begin{lemma}[Majorana sign-gauge invariance]\label{lem:gauge-invariance}
  For every $s\in\{\pm1\}^N$ and every $g$,
  \begin{equation}\label{eq:gauge-invariance}
    L_N(\beta,g)=L_N(\beta,\chi_sg).
  \end{equation}
\end{lemma}

\begin{proof}
  The family $(s_a\psi_a)_{a=1}^N$ satisfies the same Clifford relations as
  $(\psi_a)_{a=1}^N$.  By uniqueness of the irreducible complex Clifford
  representation, there exists a unitary $U_s$ such that
  \begin{equation}\label{eq:unitary-sign-flip}
    U_s\psi_aU_s^*=s_a\psi_a
    \qquad
    \text{for all }a\in[N].
  \end{equation}
  Hence $U_s\psi_IU_s^*=\chi_s(I)\psi_I$ and
  \begin{equation}\label{eq:H-gauge}
    U_sH_N(g)U_s^*=H_N(\chi_sg).
  \end{equation}
  The normalized trace is unitary invariant, so $Z_N(\beta,g)=Z_N(\beta,\chi_sg)$.
  Taking logarithms proves the claim.
\end{proof}

\begin{lemma}[Gauge parity selection rule]\label{lem:selection}
  Let $\alpha$ be a multi-index.  If there exists $a\in[N]$ such that
  \begin{equation}\label{eq:selection-rule}
    \sum_{I\ni a}\alpha_I\equiv1\pmod2,
  \end{equation}
  then $c_{N,\alpha}(\beta)=0$.  In particular,
  \begin{equation}\label{eq:first-chaos-zero}
    P_1L_N(\beta)=0,
  \end{equation}
  and for $I\neq J$,
  \begin{equation}\label{eq:off-diag-second-zero}
    \E[\partial_{g_I}\partial_{g_J}L_N(\beta)]=0.
  \end{equation}
\end{lemma}

\begin{proof}
  Since $\chi_sg$ has the same law as $g$ and $L_N(\beta,\chi_sg)=L_N(\beta,g)$,
  \begin{equation}\label{eq:coeff-gauge}
    c_{N,\alpha}(\beta)
    =
    \frac1{\alpha!}\E[L_N(\beta,g)H_\alpha(\chi_sg)].
  \end{equation}
  Since $h_k(-x)=(-1)^kh_k(x)$,
  \begin{equation}\label{eq:Hermite-sign}
    H_\alpha(\chi_sg)
    =
    \prod_{a=1}^N s_a^{\sum_{I\ni a}\alpha_I}H_\alpha(g).
  \end{equation}
  If \eqref{eq:selection-rule} holds, choose $s_a=-1$ and $s_b=1$ for
  $b\neq a$.  Then $c_{N,\alpha}(\beta)=-c_{N,\alpha}(\beta)$, so
  $c_{N,\alpha}(\beta)=0$.  The first-chaos statement follows by taking
  $\alpha=e_I$.  If $I\neq J$, choose $a\in I\triangle J$ and take
  $\alpha=e_I+e_J$.
\end{proof}

\subsection{Exact second chaos}

\begin{lemma}[First derivative identity]\label{lem:first-derivative}
  For every $I\in\calI_N$,
  \begin{equation}\label{eq:first-derivative}
    \partial_{g_I}L_N(\beta)=-\beta\sigma_N\langle\psi_I\rangle_{\beta,g}.
  \end{equation}
\end{lemma}

\begin{proof}
  By the derivative formula for the matrix exponential,
  \begin{equation}\label{eq:exp-derivative}
    \partial_{g_I}e^{-\beta H_N(g)}
    =
    -\beta\sigma_N
    \int_0^1
      e^{-(1-t)\beta H_N(g)}\psi_I e^{-t\beta H_N(g)}
    \dd t.
  \end{equation}
  After taking normalized trace, cyclicity makes the integrand independent of
  $t$ and equal to $\tr(\psi_Ie^{-\beta H_N(g)})$.  Dividing by $Z_N(\beta)$
  proves \eqref{eq:first-derivative}.
\end{proof}

\begin{lemma}[Exact second chaos]\label{lem:second-chaos}
  Set
  \begin{equation}\label{eq:an-explicit-def}
    a_N(\beta)=\frac{\beta\partial_\beta\E L_N(\beta)}{2M_N}.
  \end{equation}
  Then
  \begin{equation}\label{eq:p2-exact}
    P_2L_N(\beta)=a_N(\beta)\sum_{I\in\calI_N}(g_I^2-1).
  \end{equation}
  Moreover,
  \begin{equation}\label{eq:an-coeff}
    2M_Na_N(\beta)=\beta\partial_\beta\E L_N(\beta).
  \end{equation}
\end{lemma}

\begin{proof}
  The second chaos has the form
  \begin{equation}\label{eq:p2-general}
    P_2L_N(\beta)
    =
    \sum_{I\in\calI_N}c_{2e_I}(\beta)(g_I^2-1)
    +
    \sum_{I<J}c_{e_I+e_J}(\beta)g_Ig_J.
  \end{equation}
  By Lemma~\ref{lem:selection}, the off-diagonal coefficients vanish.  The
  diagonal coefficient is independent of $I$ by permutation symmetry of the
  Majorana labels and exchangeability of the disorder.  Thus
  \begin{equation*}
    P_2L_N(\beta)
    =
    \widetilde a_N(\beta)\sum_{I\in\calI_N}(g_I^2-1)
  \end{equation*}
  for the common diagonal coefficient $\widetilde a_N(\beta)$.

  By Lemma~\ref{lem:hermite-coeff},
  \begin{equation}\label{eq:an-second-derivative}
    \widetilde a_N(\beta)=\frac12\E[\partial_{g_I}^2L_N(\beta)].
  \end{equation}
  Summing over $I$ and using Gaussian integration by parts,
  \begin{equation}\label{eq:gibp-second-chaos}
    2M_N\widetilde a_N(\beta)
    =
    \sum_{I\in\calI_N}\E[\partial_{g_I}^2L_N(\beta)]
    =
    \sum_{I\in\calI_N}\E[g_I\partial_{g_I}L_N(\beta)].
  \end{equation}
  By Lemma~\ref{lem:first-derivative},
  \begin{equation}\label{eq:g-derivative-sum}
    \sum_{I\in\calI_N}g_I\partial_{g_I}L_N(\beta)
    =
    -\beta\sigma_N\sum_{I\in\calI_N}g_I\langle\psi_I\rangle_{\beta,g}
    =
    \beta\partial_\beta L_N(\beta).
  \end{equation}
  Taking expectation proves
  $\widetilde a_N(\beta)=\beta\partial_\beta\E L_N(\beta)/(2M_N)$, hence
  \eqref{eq:an-explicit-def} and \eqref{eq:an-coeff}.
\end{proof}

\begin{proposition}[Exact reduction to the diagonal second chaos]
  \label{prop:exact-reduction}
  For every fixed $\beta\in\R$,
  \begin{equation}\label{eq:exact-reduction}
    L_N(\beta)-\E L_N(\beta)
    =
    \frac{\beta\partial_\beta\E L_N(\beta)}{2M_N}
    \sum_{I\in\calI_N}(g_I^2-1)
    +
    R_N(\beta),
  \end{equation}
  where
  \begin{equation}\label{eq:Rn}
    R_N(\beta)=\sum_{r\geq3}P_rL_N(\beta).
  \end{equation}
\end{proposition}

\begin{proof}
  Combine the Hermite decomposition, the identity $P_1L_N=0$, and
  Lemma~\ref{lem:second-chaos}.
\end{proof}

\subsection{The second-chaos CLT}

\begin{lemma}[Asymptotic coefficient]\label{lem:coefficient-limit}
  Under Assumption~\ref{ass:local-lln}, for every compact
  $K\subset(-\beta_0,\beta_0)$,
  \begin{equation}\label{eq:coefficient-limit}
    \sup_{\beta\in K}
    \left|
      N^{q/2-1}\sqrt{2M_N}\,a_N(\beta)
      -
      \sqrt{\frac{q!}{2}}\,\beta\ell_q'(\beta)
    \right|
    \to0.
  \end{equation}
\end{lemma}

\begin{proof}
  By \eqref{eq:an-coeff},
  \begin{equation}\label{eq:an-from-derivative}
    a_N(\beta)=\frac{\beta\partial_\beta\E L_N(\beta)}{2M_N}.
  \end{equation}
  Proposition~\ref{prop:derivative-lln-from-lln} gives, uniformly for
  $\beta\in K$,
  \begin{equation}\label{eq:derivative-asymptotic}
    \partial_\beta\E L_N(\beta)=N\ell_q'(\beta)+o(N).
  \end{equation}
  Since $M_N=\binom{N}{q}\sim N^q/q!$,
  \begin{equation}\label{eq:Nn-asymptotic}
    \frac{N^{q/2}}{\sqrt{2M_N}}\to\sqrt{\frac{q!}{2}}.
  \end{equation}
  Substituting these estimates into \eqref{eq:an-from-derivative} proves the
  result.
\end{proof}

\begin{lemma}[Second-chaos CLT]\label{lem:p2-clt}
  Under Assumption~\ref{ass:local-lln}, for any fixed
  $\beta_1,\ldots,\beta_k\in(-\beta_0,\beta_0)$,
  \begin{equation}\label{eq:p2-clt}
    \left(
      N^{q/2-1}P_2L_N(\beta_1),\ldots,
      N^{q/2-1}P_2L_N(\beta_k)
    \right)
    \Rightarrow
    \left(
      \sqrt{\frac{q!}{2}}\,\beta_1\ell_q'(\beta_1)\xi,
      \ldots,
      \sqrt{\frac{q!}{2}}\,\beta_k\ell_q'(\beta_k)\xi
    \right),
  \end{equation}
  where $\xi\sim\N(0,1)$.
\end{lemma}

\begin{proof}
  Define
  \begin{equation}\label{eq:Sn}
    S_n=
    \frac1{\sqrt{2M_N}}
    \sum_{I\in\calI_N}(g_I^2-1).
  \end{equation}
  Since the variables $g_I^2-1$ are i.i.d., centered, have variance $2$, and
  have finite third moment, the Lindeberg CLT gives
  \begin{equation}\label{eq:Sn-clt}
    S_n\Rightarrow\xi,
    \qquad
    \xi\sim\N(0,1).
  \end{equation}
  By \eqref{eq:p2-exact},
  \begin{equation}\label{eq:p2-Sn}
    N^{q/2-1}P_2L_N(\beta)
    =
    \left(N^{q/2-1}\sqrt{2M_N}\,a_N(\beta)\right)S_n.
  \end{equation}
  Lemma~\ref{lem:coefficient-limit} identifies the deterministic coefficient.
  The joint convergence follows because all coordinates are deterministic
  multiples of the same scalar $S_n$.
\end{proof}

\subsection{Higher-chaos closure from two-marked susceptibility}

The next two elementary facts replace the former all-orders marked-tail
assumption by a second-derivative estimate.

\begin{lemma}[Kubo covariance identity]
  \label{lem:kubo-hessian}
  For every $I,J\in\calI_N$ and every $\beta\in\R$,
  \begin{equation}\label{eq:hessian-kubo}
    \partial_{g_I}\partial_{g_J}L_N(\beta)
    =
    \beta^2\sigma_N^2\mathcal C_{I,J}(\beta,g).
  \end{equation}
  Consequently, under Assumption~\ref{ass:two-marked-susceptibility}, for
  every compact $K\subset(-\beta_{\rm sus},\beta_{\rm sus})$, with
  $\beta_K=\sup_{\beta\in K}|\beta|$,
  \begin{equation}\label{eq:hessian-hs-bound}
    \sup_{\beta\in K}
    \E
    \sum_{I,J\in\calI_N}
    \left|
      \partial_{g_I}\partial_{g_J}L_N(\beta)
    \right|^2
    \leq
    \beta_K^4\frac{((q-1)!)^2}{q!}S_{q,K}N^{2-q}.
  \end{equation}
\end{lemma}

\begin{proof}
  Lemma~\ref{lem:first-derivative} gives
  \begin{equation}\label{eq:first-derivative-before-kubo}
    \partial_{g_I}L_N(\beta)
    =
    -\beta\sigma_N\langle\psi_I\rangle_{\beta,g}.
  \end{equation}
  Differentiate $\langle\psi_I\rangle_{\beta,g}$ with respect to $g_J$.
  Using the Duhamel formula
  \begin{equation}\label{eq:duhamel-kubo}
    \partial_{g_J}e^{-\beta H_N(g)}
    =
    -\beta\sigma_N
    \int_0^1
      e^{-s\beta H_N(g)}\psi_J e^{-(1-s)\beta H_N(g)}
    \dd s
  \end{equation}
  and differentiating the denominator $Z_N(\beta)$ gives
  \begin{align}\label{eq:derivative-thermal-average}
    \partial_{g_J}\langle\psi_I\rangle_{\beta,g}
    &=
    -\beta\sigma_N
    \int_0^1
    \left\{
      \frac{
        \tr\left(\psi_I e^{-s\beta H_N(g)}
        \psi_J e^{-(1-s)\beta H_N(g)}\right)
      }{Z_N(\beta)}
      -
      \langle\psi_I\rangle_{\beta,g}
      \langle\psi_J\rangle_{\beta,g}
    \right\}
    \dd s \\
    &=
    -\beta\sigma_N\mathcal C_{I,J}(\beta,g).
  \end{align}
  This proves \eqref{eq:hessian-kubo}.  By
  Assumption~\ref{ass:two-marked-susceptibility},
  \begin{align}\label{eq:hessian-hs-bound-proof}
    \sup_{\beta\in K}
    \E
    \sum_{I,J\in\calI_N}
    \left|
      \partial_{g_I}\partial_{g_J}L_N(\beta)
    \right|^2
    &\leq
    \beta_K^4\sigma_N^4
    \sum_{I\in\calI_N}
    S_{q,K} \\
    &\leq
    \beta_K^4S_{q,K}\sigma_N^4M_N
    \leq
    \beta_K^4\frac{((q-1)!)^2}{q!}S_{q,K}N^{2-q},
  \end{align}
  because $M_N\leq N^q/q!$ and
  $\sigma_N^4=((q-1)!)^2N^{-2(q-1)}$.
\end{proof}

\begin{lemma}[Second Malliavin derivative and chaos tail]
  \label{lem:malliavin-second-tail}
  Let $F\in\mathbb D^{2,2}(\gamma_N)$ have Hermite expansion
  $F=\sum_{r\geq0}P_rF$.  Then
  \begin{equation}\label{eq:second-malliavin-chaos}
    \E\norm{D^2F}_{\HS}^2
    =
    \sum_{r\geq2}r(r-1)\norm{P_rF}_2^2.
  \end{equation}
  In particular, for every $R\geq2$,
  \begin{equation}\label{eq:tail-from-hessian}
    \sum_{r>R}\norm{P_rF}_2^2
    \leq
    \frac{1}{R(R-1)}
    \E\norm{D^2F}_{\HS}^2.
  \end{equation}
\end{lemma}

\begin{proof}
  It is enough to check the identity for a finite Hermite polynomial and then
  pass to the limit by monotone convergence of the chaos-side and closability of
  the Malliavin derivative.  For a multi-index $\alpha$,
  \begin{equation}\label{eq:second-derivative-Hermite}
    \partial_{g_I}\partial_{g_J}H_\alpha
    =
    \begin{cases}
      \alpha_I\alpha_J
      H_{\alpha-e_I-e_J}, & I\neq J, \\
      \alpha_I(\alpha_I-1)H_{\alpha-2e_I}, & I=J.
    \end{cases}
  \end{equation}
  Using $\E[H_\alpha^2]=\alpha!$ and summing over $I,J$, one gets, for
  $|\alpha|=r$,
  \begin{equation}\label{eq:Hessian-Halpha}
    \E\norm{D^2H_\alpha}_{\HS}^2
    =
    r(r-1)\alpha!.
  \end{equation}
  Orthogonality of Hermite polynomials of different multi-index then gives
  \eqref{eq:second-malliavin-chaos}.  Since $r(r-1)\geq R(R-1)$ for
  $r>R$, \eqref{eq:tail-from-hessian} follows.
\end{proof}

\begin{proposition}[Two-marked susceptibility implies higher-chaos closure]
  \label{prop:higher-from-two-marked}
  Under Assumption~\ref{ass:two-marked-susceptibility}, for every compact
  $K\subset(-\beta_{\rm sus},\beta_{\rm sus})$,
  \begin{equation}\label{eq:higher-chaos-closure}
    \sup_{\beta\in K}
    \E\left[
      \left|
        N^{q/2-1}\sum_{r\geq3}P_rL_N(\beta)
      \right|^2
    \right]
    \to0.
  \end{equation}
\end{proposition}

\begin{proof}
  Orthogonality of homogeneous chaoses gives
  \begin{equation}\label{eq:higher-orthogonal-sum}
    \E\left[
      \left|
        N^{q/2-1}\sum_{r\geq3}P_rL_N(\beta)
      \right|^2
    \right]
    =
    N^{q-2}\sum_{r\geq3}\norm{P_rL_N(\beta)}_2^2.
  \end{equation}
  Fix $R\geq3$.  Split the sum into $3\leq r\leq R$ and $r>R$.
  Proposition~\ref{prop:finite-chaos-closure} gives
  \begin{equation}\label{eq:finite-part-vanishes}
    \sup_{\beta\in K}
    N^{q-2}
    \sum_{r=3}^R\norm{P_rL_N(\beta)}_2^2
    \to0.
  \end{equation}
  For the tail, apply Lemma~\ref{lem:malliavin-second-tail} to
  $F=L_N(\beta)-\E L_N(\beta)$.  Since $D^2F=D^2L_N(\beta)$,
  Lemma~\ref{lem:kubo-hessian} gives, with
  $\beta_K=\sup_{\beta\in K}|\beta|$,
  \begin{equation}\label{eq:tail-hessian-bound}
    \sup_{\beta\in K}
    N^{q-2}
    \sum_{r>R}\norm{P_rL_N(\beta)}_2^2
    \leq
    \frac{\beta_K^4((q-1)!)^2S_{q,K}}{q!R(R-1)}.
  \end{equation}
  Taking $\limsup_{N\to\infty}$ and then $R\to\infty$ proves
  \eqref{eq:higher-chaos-closure}.
\end{proof}

\begin{proof}[Proof of Theorem~\ref{thm:main-clt}]
  By Proposition~\ref{prop:exact-reduction},
  \begin{equation}\label{eq:Xn-decomposition}
    X_N(\beta)
    =
    N^{q/2-1}P_2L_N(\beta)
    +
    N^{q/2-1}R_N(\beta).
  \end{equation}
  Lemma~\ref{lem:p2-clt} gives the finite-dimensional convergence of the first
  term.  Proposition~\ref{prop:higher-from-two-marked} shows that the second
  term is negligible in $L^2$, uniformly on compact subintervals of
  $(-\beta_*,\beta_*)$.  Slutsky's theorem proves the finite-dimensional CLT.
\end{proof}

\subsection{Fixed finite chaos estimates}

This subsection proves the fixed finite chaos estimate used in the Hessian-tail
argument.  Thus the only infinite-tail input needed for the CLT is the
second-derivative estimate supplied by
Assumption~\ref{ass:two-marked-susceptibility}.

\begin{lemma}[Log-partition derivative bound]
  \label{lem:log-partition-derivative-bound}
  Let $K$ and $B_1,\ldots,B_r$ be Hermitian matrices.  Put
  \begin{equation}\label{eq:logZ-general}
    \Phi(x_1,\ldots,x_r)
    =
    \log\tr\exp\left(K+\sum_{j=1}^r x_jB_j\right).
  \end{equation}
  Then
  \begin{equation}\label{eq:log-partition-derivative-bound}
    \left|
      \partial_{x_1}\cdots\partial_{x_r}\Phi(0)
    \right|
    \leq
    2^r r!\prod_{j=1}^r\norm{B_j}.
  \end{equation}
\end{lemma}

\begin{proof}
  Let $\mathcal Z(x)=\tr\exp(K+\sum_jx_jB_j)$ and $Z_0=\mathcal Z(0)$.  The
  Duhamel formula gives, for distinct variables $x_{i_1},\ldots,x_{i_m}$,
  \begin{equation}\label{eq:DmZ}
    \partial_{i_1}\cdots\partial_{i_m}\mathcal Z(0)
    =
    \sum_{\pi\in S_m}
    \int_{\Delta_m}
    \tr\left(
      e^{t_0K}B_{i_{\pi(1)}}e^{t_1K}\cdots
      B_{i_{\pi(m)}}e^{t_mK}
    \right)\dd t.
  \end{equation}
  Here $\Delta_m=\{t_j\geq0:\sum_{j=0}^m t_j=1\}$.  By noncommutative Holder,
  each trace in \eqref{eq:DmZ} is bounded in absolute value by
  \begin{equation}\label{eq:holder-bound-DmZ}
    Z_0\prod_{j=1}^m\norm{B_{i_j}}.
  \end{equation}
  Therefore
  \begin{equation}\label{eq:DmZ-bound}
    \left|
      \frac{\partial_{i_1}\cdots\partial_{i_m}\mathcal Z(0)}{Z_0}
    \right|
    \leq
    m!\prod_{j=1}^m\norm{B_{i_j}}.
  \end{equation}
  The derivative of $\log\mathcal Z$ is obtained from the usual moment-cumulant
  formula over set partitions:
  \begin{equation}\label{eq:partition-log-derivative}
    \partial_1\cdots\partial_r\log\mathcal Z(0)
    =
    \sum_{\Pi\in\mathsf P([r])}
    (-1)^{|\Pi|-1}(|\Pi|-1)!
    \prod_{A\in\Pi}
    \frac{\partial_A\mathcal Z(0)}{Z_0}.
  \end{equation}
  Applying \eqref{eq:DmZ-bound} to every block gives
  \begin{equation}\label{eq:partition-bound}
    \left|
      \partial_1\cdots\partial_r\log\mathcal Z(0)
    \right|
    \leq
    \left(
      \sum_{\Pi\in\mathsf P([r])}(|\Pi|-1)!\prod_{A\in\Pi}|A|!
    \right)
    \prod_{j=1}^r\norm{B_j}.
  \end{equation}
  For partitions with $k$ blocks,
  \begin{equation}\label{eq:block-factor-count}
    \sum_{\substack{\Pi\in\mathsf P([r])\\ |\Pi|=k}}
    \prod_{A\in\Pi}|A|!
    =
    \frac{r!}{k!}\binom{r-1}{k-1}.
  \end{equation}
  Hence the combinatorial factor is
  \begin{equation}\label{eq:partition-factor-explicit}
    r!\sum_{k=1}^r\frac1k\binom{r-1}{k-1}
    =
    r!\frac{2^r-1}{r}
    \leq
    2^rr!.
  \end{equation}
  This proves \eqref{eq:log-partition-derivative-bound}.  Repeated variables
  follow by polarization or by the same formula with repeated labels.
\end{proof}

For a multi-index $\alpha$, call it even if
\begin{equation}\label{eq:alpha-even-def}
  \sum_{I\ni a}\alpha_I\equiv0\pmod2
  \qquad\text{for all }a\in[N].
\end{equation}
Only even multi-indices survive by Lemma~\ref{lem:selection}.

\begin{lemma}[Counting even marked hypergraphs at fixed order]
  \label{lem:even-marked-count}
  Fix an even integer $q$ and an integer $R\geq1$, and set
  \begin{equation}\label{eq:even-count-constant}
    E_{q,R}
    =
    \left(\left\lfloor\frac{Rq}{2}\right\rfloor+1\right)
    \left(\frac{Rq}{2}\right)^{qR}.
  \end{equation}
  Then, for every $1\leq r\leq R$ and every $N$,
  \begin{equation}\label{eq:even-marked-count}
    \sum_{\substack{|\alpha|=r\\ \alpha\text{ even}}}\frac1{\alpha!}
    \leq
    E_{q,R}N^{rq/2}.
  \end{equation}
\end{lemma}

\begin{proof}
  If $\alpha$ is even and $|\alpha|=r$, the total marked incidence degree is
  $rq$ and every vertex touched by the marked hyperedges has degree at least
  $2$.  Hence the union of the marked edges has size at most $rq/2$.  The number
  of ordered $r$-tuples $(I_1,\ldots,I_r)$ with even incidence is bounded by
  \begin{equation}\label{eq:ordered-even-count}
    \sum_{v\leq rq/2}\binom{N}{v}\binom{v}{q}^r
    \leq E_{q,R}N^{rq/2},
    \qquad 1\leq r\leq R.
  \end{equation}
  A multi-index $\alpha$ with $|\alpha|=r$ corresponds to $r!/\alpha!$
  ordered tuples.  Therefore
  \begin{equation}\label{eq:multiindex-from-ordered-count}
    \sum_{\substack{|\alpha|=r\\ \alpha\text{ even}}}\frac1{\alpha!}
    =
    \frac1{r!}
    \#\{(I_1,\ldots,I_r):\text{ the induced incidence is even}\},
  \end{equation}
  and the preceding ordered-tuple bound proves the claim.
\end{proof}

\begin{proposition}[Finite-chaos closure]\label{prop:finite-chaos-closure}
  Fix an integer $R\geq3$ and a compact $K\subset\R$, and put
  $\beta_K=\sup_{\beta\in K}|\beta|$.  Then
  \begin{align}\label{eq:finite-chaos-explicit-bound}
    \sup_{\beta\in K}
    \sum_{r=3}^R\norm{P_rL_N(\beta)}_2^2
    &\leq
    E_{q,R}
    \sum_{r=3}^R
    (2^rr!)^2\beta_K^{2r}((q-1)!)^r
    N^{r(1-q/2)}.
  \end{align}
  In particular,
  \begin{equation}\label{eq:finite-chaos-closure}
    \sup_{\beta\in K}
    \sum_{r=3}^R\norm{P_rL_N(\beta)}_2^2
    =
    o(N^{2-q}).
  \end{equation}
\end{proposition}

\begin{proof}
  For $|\alpha|=r$, write $\alpha=e_{I_1}+\cdots+e_{I_r}$ with repetitions.
  Applying Lemma~\ref{lem:log-partition-derivative-bound} with
  $B_j=-\beta\sigma_N\psi_{I_j}$ gives, for $\beta\in K$ and $r\leq R$,
  \begin{equation}\label{eq:fixed-alpha-bound}
    \left|\E\partial^\alpha L_N(\beta)\right|
    \leq
    2^rr!\beta_K^r\sigma_N^r.
  \end{equation}
  By Lemma~\ref{lem:hermite-coeff} and orthogonality,
  \begin{equation}\label{eq:Pr-norm-coeff}
    \norm{P_rL_N(\beta)}_2^2
    =
    \sum_{|\alpha|=r}\frac1{\alpha!}
    \left(\E\partial^\alpha L_N(\beta)\right)^2.
  \end{equation}
  Lemma~\ref{lem:selection} restricts the sum to even $\alpha$.  Hence, using
  Lemma~\ref{lem:even-marked-count},
  \begin{equation}\label{eq:Pr-fixed-bound}
    \norm{P_rL_N(\beta)}_2^2
    \leq
    E_{q,R}(2^rr!)^2\beta_K^{2r}
    \sigma_N^{2r}N^{rq/2}
    =
    E_{q,R}(2^rr!)^2\beta_K^{2r}((q-1)!)^r
    N^{r(1-q/2)}.
  \end{equation}
  Summing \eqref{eq:Pr-fixed-bound} gives
  \eqref{eq:finite-chaos-explicit-bound}.  For $q\geq4$ and $r\geq3$,
  $r(1-q/2)<2-q$, so the explicit bound implies
  \eqref{eq:finite-chaos-closure}.
\end{proof}



\section{Proof attempt for the two-marked susceptibility estimate}
\label{sec:two-marked-proof-attempt}

This section pushes the proof of
Assumption~\ref{ass:two-marked-susceptibility} to a precise high-temperature
cluster-expansion input.  The source-derivative identity, the finite-order Wick
reductions, the Duhamel cumulant tree bound, and the connected parity counting
lemma below are rigorous.  These results explain why the desired estimate has
the correct power of $N$ and why the remaining small parameter should be
$\beta^2$.  They do not, however, by themselves prove the needed summability:
repeated hyperedges make a coefficientwise or occurrence-level tree summation
too crude.  The only unproved input left in this section is therefore stated as
a separate resummed two-replica $L^2$ cluster-expansion estimate,
Assumption~\ref{ass:two-replica-cluster-expansion}.

At $\beta=0$, the Kubo covariance is diagonal.  Indeed,
Lemma~\ref{lem:trace-orth} below gives
\begin{equation}\label{eq:kubo-beta-zero}
  \mathcal C_{I,J}(0,g)
  =
  \tr(\psi_I\psi_J)
  =
  \ind_{I=J}.
\end{equation}
Hence
\begin{equation}\label{eq:kubo-beta-zero-sum}
  \sup_{I\in\calI_N}
  \sum_{J\in\calI_N}
  \left|\mathcal C_{I,J}(0,g)\right|^2
  =
  1.
\end{equation}
The target estimate says that this bounded row-susceptibility persists, in
annealed $L^2$, for sufficiently small $|\beta|$.

\subsection{Source derivatives and finite-order connected coefficients}

For fixed $I,J\in\calI_N$ introduce source variables $u,v$ and set
\begin{equation}\label{eq:source-log-partition}
  \Lambda_{I,J}(u,v;\beta,g)
  =
  \log\tr\exp\left(-\beta H_N(g)+u\psi_I+v\psi_J\right).
\end{equation}

\begin{lemma}[Kubo covariance as a source derivative]
  \label{lem:kubo-source-derivative}
  For every $I,J\in\calI_N$,
  \begin{equation}\label{eq:kubo-source-derivative}
    \mathcal C_{I,J}(\beta,g)
    =
    \partial_u\partial_v\Lambda_{I,J}(u,v;\beta,g)\big|_{u=v=0}.
  \end{equation}
\end{lemma}

\begin{proof}
  First differentiate \eqref{eq:source-log-partition} with respect to $u$.
  By the Duhamel formula and cyclicity of trace,
  \begin{equation}\label{eq:first-source-derivative}
    \partial_u\Lambda_{I,J}(u,v;\beta,g)\big|_{u=v=0}
    =
    \langle\psi_I\rangle_{\beta,g}.
  \end{equation}
  Differentiating this identity with respect to $v$ and applying the Duhamel
  formula once more gives
  \begin{align}\label{eq:second-source-derivative-proof}
    \partial_u\partial_v\Lambda_{I,J}(u,v;\beta,g)\big|_{u=v=0}
    &=
    \int_0^1
    \left\{
      \frac{
        \tr\left(\psi_I e^{-s\beta H_N(g)}
        \psi_J e^{-(1-s)\beta H_N(g)}\right)
      }{Z_N(\beta)}
      -
      \langle\psi_I\rangle_{\beta,g}
      \langle\psi_J\rangle_{\beta,g}
    \right\}
    \dd s.
  \end{align}
  This is exactly \eqref{eq:kubo-covariance}.
\end{proof}

For $m\geq0$ and $K_1,\ldots,K_m\in\calI_N$, define the connected Duhamel
coefficient
\begin{equation}\label{eq:connected-duhamel-coeff}
  \kappa(I,J;K_1,
  \ldots,K_m)
  =
  \partial_u\partial_v\partial_{t_1}\cdots\partial_{t_m}
  \log\tr\exp\left(
    u\psi_I+v\psi_J+
    \sum_{a=1}^m t_a\psi_{K_a}
  \right)\bigg|_{u=v=t_1=\cdots=t_m=0}.
\end{equation}
For $M\geq0$, let the $M$-th formal truncation be
\begin{equation}\label{eq:truncated-connected-expansion}
  \mathcal C^{[M]}_{I,J}(\beta,g)
  =
  \sum_{m=0}^M
  \frac{(-\beta\sigma_N)^m}{m!}
  \sum_{K_1,
  \ldots,K_m\in\calI_N}
  g_{K_1}\cdots g_{K_m}
  \kappa(I,J;K_1,
  \ldots,K_m).
\end{equation}

The word ``connected'' in \eqref{eq:connected-duhamel-coeff} is literal.  If the
intersection hypergraph on the hyperedges $I,J,K_1,\ldots,K_m$ splits into two
vertex-disjoint subfamilies separating $I$ from $J$, then
\begin{equation}\label{eq:connected-vanishing}
  \kappa(I,J;K_1,
  \ldots,K_m)=0.
\end{equation}
Indeed, normalized trace factorizes over disjoint Majorana tensor factors, and
the logarithm extracts connected components by the exponential formula.

\subsection{Cumulant structure of the Duhamel coefficients}
\label{subsec:duhamel-cumulant-structure}

We next isolate what can be proved directly about the coefficients
\eqref{eq:connected-duhamel-coeff}.  This is the useful replacement for an
informal statement that ``the logarithm keeps only connected diagrams.''  In
particular, the sharp dependency-graph tree bound for the zero-source
Duhamel coefficients is proved below.

Let $r\geq1$ and let $E_1,\ldots,E_r\in\calI_N$.  Put
$A_a=\psi_{E_a}$.  For $S\subset[r]$, define the symmetrized trace moment
\begin{equation}\label{eq:sym-trace-moment}
  \mu_{\mathbf E}(S)
  =
  \partial_{(t_a)_{a\in S}}
  \tr\exp\left(\sum_{a\in S}t_aA_a\right)\bigg|_{t=0}.
\end{equation}
Equivalently, if $S=\{a_1,\ldots,a_s\}$, then
\begin{equation}\label{eq:sym-trace-moment-permutation}
  \mu_{\mathbf E}(S)
  =
  \frac1{s!}
  \sum_{\pi\in\mathfrak S_s}
  \tr\left(A_{a_{\pi(1)}}\cdots A_{a_{\pi(s)}}\right),
  \qquad
  \mu_{\mathbf E}(\emptyset)=1.
\end{equation}
Let $G_{\mathbf E}$ be the intersection graph on $[r]$, where $a\sim b$ iff
$E_a\cap E_b\neq\emptyset$.

\begin{lemma}[Duhamel coefficient as an ordinary cumulant]
  \label{lem:duhamel-as-cumulant}
  For $E_1,\ldots,E_r\in\calI_N$,
  \begin{equation}\label{eq:duhamel-as-cumulant}
    \partial_{t_1}\cdots\partial_{t_r}
    \log\tr\exp\left(\sum_{a=1}^r t_a\psi_{E_a}\right)\bigg|_{t=0}
    =
    \sum_{\pi\in\Pi_r}
    (|\pi|-1)!(-1)^{|\pi|-1}
    \prod_{B\in\pi}\mu_{\mathbf E}(B),
  \end{equation}
  where $\Pi_r$ denotes the set of set partitions of $[r]$.
\end{lemma}

\begin{proof}
  Set
  \begin{equation}\label{eq:moment-generating-function}
    M(t_1,\ldots,t_r)
    =
    \tr\exp\left(\sum_{a=1}^r t_a\psi_{E_a}\right).
  \end{equation}
  Since $M(0)=1$, the Taylor coefficient of $\log M$ is obtained from the
  standard formal identity
  \begin{equation}\label{eq:log-moment-cumulant-formal}
    \log M
    =
    \sum_{k\geq1}\frac{(-1)^{k-1}}{k}(M-1)^k.
  \end{equation}
  Taking the coefficient of $t_1\cdots t_r$ in
  \eqref{eq:log-moment-cumulant-formal} gives the partition formula
  \eqref{eq:duhamel-as-cumulant}.  The moments that enter are precisely
  \eqref{eq:sym-trace-moment}, because differentiating the exponential at the
  origin gives the symmetrized trace expression
  \eqref{eq:sym-trace-moment-permutation}.
\end{proof}

\begin{lemma}[Factorization and parity]
  \label{lem:moment-factorization-parity}
  The moments $\mu_{\mathbf E}$ satisfy the following two properties.

  First, if $S\subset[r]$ and $C_1,\ldots,C_\ell$ are the connected components
  of $G_{\mathbf E}[S]$, then
  \begin{equation}\label{eq:moment-factorization-components}
    \mu_{\mathbf E}(S)
    =
    \prod_{a=1}^\ell \mu_{\mathbf E}(C_a).
  \end{equation}
  Second,
  \begin{equation}\label{eq:moment-parity-selection}
    \mu_{\mathbf E}(S)=0
    \qquad\text{unless}\qquad
    \triangle_{a\in S}E_a=\emptyset,
  \end{equation}
  where $\triangle$ denotes symmetric difference.
  Consequently,
  \begin{equation}\label{eq:duhamel-coeff-disconnected-or-odd-zero}
    \partial_{t_1}\cdots\partial_{t_r}
    \log\tr\exp\left(\sum_{a=1}^r t_a\psi_{E_a}\right)\bigg|_{t=0}=0
  \end{equation}
  whenever $G_{\mathbf E}$ is disconnected or
  $\triangle_{a=1}^rE_a\neq\emptyset$.
\end{lemma}

\begin{proof}
  If $C_1,\ldots,C_\ell$ are the connected components of $G_{\mathbf E}[S]$,
  then the Majorana labels used by different $C_a$ are disjoint.  The Clifford
  algebra representation decomposes as a tensor product over these disjoint
  label sets, and the normalized trace factorizes.  Since monomials supported
  on different components commute, the symmetrized product in
  \eqref{eq:sym-trace-moment-permutation} factorizes as well.  This proves
  \eqref{eq:moment-factorization-components}.

  For the parity statement, suppose some Majorana label $b$ belongs to an odd
  number of sets $E_a$, $a\in S$.  Flip the sign of $\psi_b$ and leave all other
  Majorana generators fixed.  This is a Clifford automorphism and hence is
  unitarily implemented.  Therefore normalized traces are invariant under this
  operation.  On the other hand every product appearing in
  \eqref{eq:sym-trace-moment-permutation} changes sign.  Hence
  $\mu_{\mathbf E}(S)=-\mu_{\mathbf E}(S)$, proving
  \eqref{eq:moment-parity-selection}.

  Finally, if $G_{\mathbf E}$ is disconnected, then
  \eqref{eq:moment-factorization-components} says that the moment functional
  factorizes across a nontrivial partition of $[r]$.  The cumulant formula
  \eqref{eq:duhamel-as-cumulant} then vanishes by the usual cancellation of
  cumulants for a product moment functional.  If
  $\triangle_{a=1}^rE_a\neq\emptyset$, then every block product in
  \eqref{eq:duhamel-as-cumulant} has at least one block with odd total
  incidence, so it vanishes by \eqref{eq:moment-parity-selection}.
\end{proof}

\begin{lemma}[A crude absolute cumulant bound]
  \label{lem:crude-cumulant-bound}
  For all $E_1,\ldots,E_r\in\calI_N$,
  \begin{equation}\label{eq:crude-cumulant-bound}
    \left|
      \partial_{t_1}\cdots\partial_{t_r}
      \log\tr\exp\left(\sum_{a=1}^r t_a\psi_{E_a}\right)\bigg|_{t=0}
    \right|
    \leq
    e^r r!.
  \end{equation}
\end{lemma}

\begin{proof}
  Since $\norm{\psi_E}\leq1$ for every $E\in\calI_N$, each symmetrized trace
  moment satisfies $|\mu_{\mathbf E}(S)|\leq1$.  Hence
  \eqref{eq:duhamel-as-cumulant} gives
  \begin{equation}\label{eq:cumulant-partition-sum-bound}
    |\kappa_r|
    \leq
    \sum_{\pi\in\Pi_r}(|\pi|-1)!.
  \end{equation}
  The number of partitions of $[r]$ into $k$ blocks is at most $k^r/k!$.
  Therefore
  \begin{equation}\label{eq:cumulant-partition-sum-bound-2}
    \sum_{\pi\in\Pi_r}(|\pi|-1)!
    \leq
    \sum_{k=1}^r \frac{k^r}{k!}(k-1)!
    \leq
    \sum_{k=1}^r k^{r-1}
    \leq
    r^r
    \leq
    e^r r!,
  \end{equation}
  where the last inequality follows from Stirling's lower bound.  This proves
  \eqref{eq:crude-cumulant-bound}.
\end{proof}

The preceding lemmas are rigorous and show exactly what the connected Duhamel
coefficient knows: it has a dependency graph, it vanishes off connected
parity-even families, and it has at most factorial growth.  We now prove the
sharper tree-graph estimate that was previously left as a hypothesis.

\begin{lemma}[Chromatic derivative bound]
  \label{lem:chromatic-derivative-bound}
  Let $G$ be a finite graph.  If $G$ is connected, then
  \begin{equation}\label{eq:chromatic-derivative-bound}
    |\chi_G'(0)|
    \leq
    \tau(G),
  \end{equation}
  where $\chi_G$ is the chromatic polynomial and $\tau(G)$ is the number of
  spanning trees of $G$.  If $G$ is disconnected, then $\chi_G'(0)=0$.
\end{lemma}

\begin{proof}
  If $G$ is disconnected, then $\chi_G(z)$ is divisible by $z^2$, and hence
  $\chi_G'(0)=0$.  Assume that $G$ is connected.  Fix an ordering of the edge
  set.  Whitney's broken-circuit theorem writes the chromatic polynomial as
  \begin{equation}\label{eq:whitney-broken-circuit}
    \chi_G(z)
    =
    \sum_{A}(-1)^{|A|}z^{c(A)},
  \end{equation}
  where the sum is over edge subsets $A$ containing no broken circuit and
  $c(A)$ is the number of connected components of the spanning subgraph
  $([r],A)$.  The coefficient of $z$ therefore counts, up to sign, the
  no-broken-circuit edge sets $A$ for which $([r],A)$ is connected.  Such an
  $A$ is necessarily a spanning tree: if it were connected and contained a
  cycle, the largest edge in that cycle would produce a broken circuit.  Hence
  $|\chi_G'(0)|$ is bounded by the number of spanning trees of $G$.
\end{proof}

\begin{lemma}[Abstract dependency-graph cumulant tree bound]
  \label{lem:abstract-dependency-tree-bound}
  Let $G$ be a graph on $[r]$.  Let $\mu:2^{[r]}\to\C$ satisfy
  $\mu(\emptyset)=1$, $|\mu(S)|\leq1$ for all $S\subset[r]$, and
  \begin{equation}\label{eq:abstract-factorization}
    \mu(S)
    =
    \prod_{C\in\Comp(G[S])}\mu(C)
  \end{equation}
  for every $S\subset[r]$.  Define
  \begin{equation}\label{eq:abstract-cumulant}
    \kappa_\mu([r])
    =
    \sum_{\pi\in\Pi_r}
    (-1)^{|\pi|-1}(|\pi|-1)!
    \prod_{B\in\pi}\mu(B).
  \end{equation}
  Then
  \begin{equation}\label{eq:abstract-tree-bound}
    |\kappa_\mu([r])|
    \leq
    2^{r-1}\tau(G),
  \end{equation}
  where $\tau(G)$ is the number of spanning trees of $G$, and $\tau(G)=0$ if
  $G$ is disconnected.
\end{lemma}

\begin{proof}
  For a partition $\pi$ of $[r]$, let $\pi^G$ be the partition obtained by
  replacing each block $B\in\pi$ by the connected components of $G[B]$.
  By \eqref{eq:abstract-factorization},
  \begin{equation}\label{eq:regroup-factorization}
    \prod_{B\in\pi}\mu(B)
    =
    \prod_{C\in\pi^G}\mu(C).
  \end{equation}
  Hence
  \begin{equation}\label{eq:cumulant-regrouped}
    \kappa_\mu([r])
    =
    \sum_{\rho}
    a_G(\rho)
    \prod_{C\in\rho}\mu(C),
  \end{equation}
  where the sum is over partitions $\rho$ whose blocks are connected in $G$, and
  \begin{equation}\label{eq:aG-rho-def}
    a_G(\rho)
    =
    \sum_{\pi:\pi^G=\rho}
    (-1)^{|\pi|-1}(|\pi|-1)!.
  \end{equation}

  Let $\overline G_\rho$ be the simple quotient graph whose vertices are the
  blocks of $\rho$, with an edge between two blocks if $G$ has at least one edge
  between them.  A partition $\pi$ satisfying $\pi^G=\rho$ is the same as a
  partition of the vertices of $\overline G_\rho$ into nonempty independent
  sets.  If $S_{\overline G_\rho}(k)$ denotes the number of such partitions into
  $k$ independent sets, then
  \begin{equation}\label{eq:aG-rho-independent}
    a_G(\rho)
    =
    \sum_{k\geq1}
    (-1)^{k-1}(k-1)!S_{\overline G_\rho}(k).
  \end{equation}
  The chromatic polynomial has the falling-factorial expansion
  \begin{equation}\label{eq:chromatic-falling-factorial}
    \chi_{\overline G_\rho}(z)
    =
    \sum_{k\geq1}
    S_{\overline G_\rho}(k)(z)_k,
    \qquad
    (z)_k=z(z-1)\cdots(z-k+1).
  \end{equation}
  Since $(z)_k'|_{z=0}=(-1)^{k-1}(k-1)!$, we get
  \begin{equation}\label{eq:aG-rho-chromatic}
    a_G(\rho)=\chi_{\overline G_\rho}'(0).
  \end{equation}
  By Lemma~\ref{lem:chromatic-derivative-bound},
  \begin{equation}\label{eq:chromatic-derivative-tree}
    |a_G(\rho)|
    =
    |\chi_{\overline G_\rho}'(0)|
    \leq
    \tau(\overline G_\rho),
  \end{equation}
  with the convention that $\tau(\overline G_\rho)=0$ when the quotient graph is
  disconnected.

  Since $|\mu(C)|\leq1$, \eqref{eq:cumulant-regrouped} and
  \eqref{eq:chromatic-derivative-tree} give
  \begin{equation}\label{eq:cumulant-tree-sum-rho}
    |\kappa_\mu([r])|
    \leq
    \sum_{\rho}\tau(\overline G_\rho),
  \end{equation}
  where only $\rho$ with connected quotient contribute.
  It remains to bound this sum by $2^{r-1}\tau(G)$.

  For each connected subset $C\subset[r]$, fix once and for all a spanning tree
  $T_C$ of $G[C]$.  For each pair of disjoint subsets $C,D$ with at least one
  edge between them, fix one canonical edge of $G$ joining $C$ to $D$.
  Given a partition $\rho$ with connected blocks and a spanning tree
  $\overline T$ of $\overline G_\rho$, form a spanning tree of $G$ by taking the
  union of all $T_C$, $C\in\rho$, and the canonical inter-block edges
  corresponding to the edges of $\overline T$.  Mark the inter-block edges in
  this tree.  This maps each pair $(\rho,\overline T)$ injectively into a pair
  $(T,F)$, where $T$ is a spanning tree of $G$ and $F\subset E(T)$.  Indeed,
  removing $F$ from $T$ recovers the blocks of $\rho$, and $F$ recovers
  $\overline T$.  Therefore
  \begin{equation}\label{eq:rho-tree-sum-bound}
    \sum_\rho \tau(\overline G_\rho)
    \leq
    \#\{(T,F):T\in\mathsf{Tree}(G),\ F\subset E(T)\}
    =
    2^{r-1}\tau(G).
  \end{equation}
  Combining \eqref{eq:cumulant-tree-sum-rho} and
  \eqref{eq:rho-tree-sum-bound} proves the lemma.
\end{proof}

\begin{proposition}[Duhamel cumulant tree-graph inequality]
  \label{prop:duhamel-tree-graph}
  For every $r\geq2$ and every $E_1,\ldots,E_r\in\calI_N$,
  \begin{equation}\label{eq:duhamel-tree-graph}
    \left|
      \partial_{t_1}\cdots\partial_{t_r}
      \log\tr\exp\left(\sum_{a=1}^r t_a\psi_{E_a}\right)\bigg|_{t=0}
    \right|
    \leq
    2^{r-1}
    \sum_{T\in\mathsf{Tree}([r])}
    \prod_{\{a,b\}\in E(T)}
    \ind_{E_a\cap E_b\neq\emptyset}.
  \end{equation}
\end{proposition}

\begin{proof}
  Let $G=G_{\mathbf E}$ be the intersection graph.  Since
  $\norm{\psi_E}\leq1$, the symmetrized trace moments satisfy
  $|\mu_{\mathbf E}(S)|\leq1$.  By
  Lemma~\ref{lem:moment-factorization-parity}, they factorize over connected
  components of $G[S]$.  Lemma~\ref{lem:duhamel-as-cumulant} identifies the
  left-hand side of \eqref{eq:duhamel-tree-graph} with the ordinary cumulant of
  this moment functional.  Lemma~\ref{lem:abstract-dependency-tree-bound}
  therefore gives
  \begin{equation}\label{eq:duhamel-tree-proof-final}
    |\kappa_D(E_1,\ldots,E_r)|
    \leq
    2^{r-1}\tau(G).
  \end{equation}
  Finally,
  \begin{equation}\label{eq:tau-expanded-tree-sum}
    \tau(G)
    =
    \sum_{T\in\mathsf{Tree}([r])}
    \prod_{\{a,b\}\in E(T)}
    \ind_{E_a\cap E_b\neq\emptyset}.
  \end{equation}
  Thus \eqref{eq:duhamel-tree-graph} holds.
\end{proof}

\subsection{The finite two-replica Wick reduction}

Squaring \eqref{eq:truncated-connected-expansion} and taking Gaussian
expectation produces two copies of the connected expansion.  The Gaussian
expectation is computed by Wick pairing.  It is useful to package the result as
diagrams.

A two-replica diagram $\Gamma$ with marks $I,J$ consists of two ordered lists of
unmarked hyperedges,
\begin{equation}\label{eq:two-replica-lists}
  \mathbf K=(K_1,
  \ldots,K_m),
  \qquad
  \mathbf L=(L_1,
  \ldots,L_{m'}),
\end{equation}
paired by a Wick pairing of the multiset
$\{K_1,
\ldots,K_m,L_1,
\ldots,L_{m'}\}$.  In particular $m+m'=2p$ is even.  Its value is
\begin{equation}\label{eq:diagram-value}
  W(\Gamma)
  =
  \frac{
    \kappa(I,J;K_1,
    \ldots,K_m)
    \kappa(I,J;L_1,
    \ldots,L_{m'})
  }{m!m'!}.
\end{equation}
Let $\mathfrak D_p^{[M]}(I,J)$ denote the set of such diagrams with
$m,m'\leq M$ and with $p$ Wick pairs for which both connected Duhamel
coefficients in \eqref{eq:diagram-value} are nonzero.

\begin{lemma}[Finite annealed two-replica formula]
  \label{lem:finite-two-replica-formula}
  For fixed $M,N,I,J$ and $\beta\in\R$,
  \begin{equation}\label{eq:finite-two-replica-formula}
    \E\left[
      \left|
        \mathcal C^{[M]}_{I,J}(\beta,g)
      \right|^2
    \right]
    =
    \sum_{p=0}^M
    (\beta\sigma_N)^{2p}
    \sum_{\Gamma\in\mathfrak D_p^{[M]}(I,J)}
    W(\Gamma).
  \end{equation}
  Consequently,
  \begin{equation}\label{eq:absolute-finite-two-replica-formula}
    \E\left[
      \left|
        \mathcal C^{[M]}_{I,J}(\beta,g)
      \right|^2
    \right]
    \leq
    \sum_{p=0}^M
    |\beta|^{2p}\sigma_N^{2p}
    \sum_{\Gamma\in\mathfrak D_p^{[M]}(I,J)}
    |W(\Gamma)|.
  \end{equation}
\end{lemma}

\begin{proof}
  The coefficients $\kappa(I,J;K_1,\ldots,K_m)$ are real.  Indeed, they are
  Taylor coefficients at the origin of the real analytic function
  \begin{equation}\label{eq:kappa-real-valued}
    (u,v,\lambda_1,\ldots,\lambda_m)
    \mapsto
    \log\tr\exp\left(
      u\psi_I+v\psi_J+\sum_{a=1}^m\lambda_a\psi_{K_a}
    \right)
  \end{equation}
  restricted to real sources.  Thus the square on the left-hand side is the
  ordinary product of the two real finite expansions.

  Insert \eqref{eq:truncated-connected-expansion} twice.  For two ordered lists
  $\mathbf K$ and $\mathbf L$, the Gaussian expectation
  \begin{equation}\label{eq:wick-expectation}
    \E\left[g_{K_1}\cdots g_{K_m}g_{L_1}\cdots g_{L_{m'}}\right]
  \end{equation}
  is zero unless $m+m'$ is even.  If $m+m'=2p$, Wick's theorem writes
  \eqref{eq:wick-expectation} as the sum over all pairings of the multiset of
  labels.  Grouping each Wick pairing with the two ordered lists gives
  $\mathfrak D_p^{[M]}(I,J)$ and the value \eqref{eq:diagram-value}.  This
  proves \eqref{eq:finite-two-replica-formula}.  The absolute bound follows by
  the triangle inequality.

\end{proof}

\subsection{Connected parity counting}

The next lemma is the purely combinatorial count that explains why the parity
constraint determining $J$ removes exactly one extra hyperedge degree of
freedom.  It is useful in any attempt to prove the cluster-expansion estimate below.

\begin{lemma}[Connected parity counting]
  \label{lem:connected-parity-counting}
  Fix $I\in\calI_N$, $p\geq0$, and $C\subset[p]$.  For
  $H_1,\ldots,H_p\in\calI_N$, set
  \begin{equation}\label{eq:J-from-cross-pairs-counting}
    J
    =
    I\triangle\bigtriangleup_{a\in C}H_a.
  \end{equation}
  Let $N_p(I,C)$ be the number of ordered $p$-tuples
  $(H_1,\ldots,H_p)$ such that $J\in\calI_N$ and the intersection graph on the
  labelled hyperedges
  \begin{equation}\label{eq:counting-labelled-hyperedges}
    I,
    \qquad
    J,
    \qquad
    H_1,\ldots,H_p
  \end{equation}
  is connected.  Then
  \begin{equation}\label{eq:connected-parity-counting}
    N_p(I,C)
    \leq
    (2e^2q^2)^p p! N^{(q-1)p}.
  \end{equation}
\end{lemma}

\begin{proof}
  The case $p=0$ is immediate.  Assume $p\geq1$ and consider a tuple counted by
  $N_p(I,C)$.
  We first prove
  \begin{equation}\label{eq:outside-union-bound-counting}
    \left|
      \left(\bigcup_{a=1}^p H_a\right)\setminus I
    \right|
    \leq
    (q-1)p.
  \end{equation}

  Consider the intersection graph on the labelled hyperedges
  $I,H_1,\ldots,H_p$, without $J$, and let its connected components be
  $\mathcal D_1,\ldots,\mathcal D_\ell$.  A hyperedge in such a component is
  called selected if it is either $I$ or $H_a$ with $a\in C$.  For a component
  $\mathcal D$, let $e(\mathcal D)$ be the number of its hyperedges and let
  $c(\mathcal D)$ be the number of its selected hyperedges.  Define
  \begin{equation}\label{eq:R-D-def}
    R(\mathcal D)
    =
    \bigtriangleup_{E\in\mathcal D,\ E\text{ selected}}E.
  \end{equation}
  Since distinct components have disjoint supports, the sets
  $R(\mathcal D_j)$ are disjoint and
  \begin{equation}\label{eq:J-disjoint-union-RD}
    J
    =
    \bigsqcup_{j=1}^\ell R(\mathcal D_j).
  \end{equation}
  Hence
  \begin{equation}\label{eq:sum-RD-q}
    \sum_{j=1}^\ell |R(\mathcal D_j)|=|J|=q.
  \end{equation}

  Each component $\mathcal D_j$ contains at least one selected hyperedge.  If it
  did not, then $R(\mathcal D_j)=\emptyset$, and the support of
  $\mathcal D_j$ would be disjoint from $J$ and from all other components.
  This would contradict the connectedness of the intersection graph on
  $I,J,H_1,\ldots,H_p$.  Thus $c(\mathcal D_j)\geq1$ for all $j$.

  Fix a component $\mathcal D$.  Let $U_{\rm sel}(\mathcal D)$ be the union of
  its selected hyperedges.  Since the total selected incidence is
  $q c(\mathcal D)$, every point in $R(\mathcal D)$ contributes at least one
  incidence and every point in
  $U_{\rm sel}(\mathcal D)\setminus R(\mathcal D)$ contributes at least two,
  so
  \begin{equation}\label{eq:selected-union-bound}
    |U_{\rm sel}(\mathcal D)|
    \leq
    \frac{q c(\mathcal D)+|R(\mathcal D)|}{2}.
  \end{equation}
  The non-selected hyperedges in $\mathcal D$ can be ordered so that each new
  one intersects the union of the previously exposed selected and non-selected
  hyperedges.  Each such hyperedge adds at most $q-1$ new vertices.  Therefore
  \begin{align}\label{eq:component-union-bound-counting}
    \left|\bigcup_{E\in\mathcal D}E\right|
    &\leq
    \frac{q c(\mathcal D)+|R(\mathcal D)|}{2}
    +(q-1)(e(\mathcal D)-c(\mathcal D)) \\
    &=
    (q-1)e(\mathcal D)
    -\frac{q-2}{2}c(\mathcal D)
    +\frac12|R(\mathcal D)|.
  \end{align}
  Summing over all components and using
  $\sum_j e(\mathcal D_j)=p+1$, $\sum_j c(\mathcal D_j)=|C|+1$, and
  \eqref{eq:sum-RD-q}, we obtain
  \begin{align}\label{eq:total-union-bound-counting}
    \left|I\cup\bigcup_{a=1}^pH_a\right|
    &\leq
    (q-1)(p+1)
    -\frac{q-2}{2}(|C|+1)
    +\frac q2 \\
    &\leq
    q+(q-1)p.
  \end{align}
  Since $|I|=q$, this proves \eqref{eq:outside-union-bound-counting}.

  We now count.  Encode $(H_1,\ldots,H_p)$ by writing each $H_a$ in increasing
  order, giving $pq$ positions.  By
  \eqref{eq:outside-union-bound-counting}, at most $(q-1)p$ distinct labels
  outside $I$ appear.  Choose one representative position for each such outside
  label and, if necessary, add arbitrary positions until exactly $(q-1)p$
  representative positions have been chosen.  The representative positions can
  be chosen in at most
  \begin{equation}\label{eq:representative-position-count}
    \binom{pq}{(q-1)p}
    =
    \binom{pq}{p}
    \leq
    (eq)^p
  \end{equation}
  ways.  Their labels can be chosen in at most $N^{(q-1)p}$ ways.  Each of the
  remaining $p$ positions is either one of the $q$ labels in $I$ or copies one
  of the representative outside labels, hence has at most $q+(q-1)p\leq qp+q$
  choices.  Thus
  \begin{equation}\label{eq:counting-final-before-factorial}
    N_p(I,C)
    \leq
    (eq)^p(qp+q)^p N^{(q-1)p}.
  \end{equation}
  Since $p+1\leq2p$ and $p^p\leq e^p p!$ for $p\geq1$, this is bounded by
  $(2e^2q^2)^p p! N^{(q-1)p}$.
\end{proof}

The finite two-replica formula is exact for every finite truncation.  It is also
useful as a diagnostic for what remains unproved.  If one bounds the two
zero-source Duhamel cumulants by occurrence-level trees and then sums the trees
absolutely after Wick pairing, repeated hyperedges create a superexponential
entropy loss.  For example, a configuration in which all $2p$ unmarked
occurrences are the same hyperedge has an occurrence-level complete graph and
therefore about $(2p)^{2p}$ spanning trees, whereas the Duhamel and Wick
factorials do not reduce this to a purely exponential factor.  This does not
show that the two-marked susceptibility is false; it shows that the
coefficientwise tree-counting strategy is too crude.  The remaining input must
therefore be a resummed two-replica $L^2$ cluster expansion, where repeated
hyperedges are treated before taking absolute values of the final cluster
weights.


\subsection{Why the remaining estimate is not proved by the naive expansions}
\label{subsec:why-naive-cluster-fails}

The preceding lemmas prove the two structural facts one expects from a
high-temperature expansion: connectedness and the correct parity count.  It is
natural to try to finish the proof by taking absolute values term by term in a
single-replica expansion, or by applying the Duhamel tree bound coefficient by
coefficient.  The following two observations explain why these two routes do
not close.  They do not disprove the resummed cluster expansion; rather, they
show why the last input must be a genuinely annealed two-replica estimate.

\begin{lemma}[Single-replica absolute activities are too large]
  \label{lem:single-replica-activity-fails}
  Let
  \begin{equation}\label{eq:single-activity-def}
    \eta_K(g)=\exp\left(c|\beta|\sigma_N|g_K|\right)-1
  \end{equation}
  for a fixed constant $c>0$.  For every fixed $I\in\calI_N$ and every fixed
  $\beta\neq0$,
  \begin{equation}\label{eq:single-activity-diverges}
    \sum_{K\in\calI_N:K\cap I\neq\emptyset}\E\eta_K(g)
    \geq
    \frac{cq|\beta|}{2}
    \sqrt{\frac{2}{\pi(q-1)!}}\,
    N^{(q-1)/2}
  \end{equation}
  for all large $N$.
\end{lemma}

\begin{proof}
  Since $e^x-1\geq x$ for $x\geq0$ and $\E|g_K|=\sqrt{2/\pi}$,
  \begin{equation}\label{eq:eta-first-order}
    \E\eta_K(g)
    \geq
    c|\beta|\sigma_N\sqrt{\frac2\pi}.
  \end{equation}
  On the other hand,
  \begin{equation}\label{eq:number-touching-I}
    \#\{K\in\calI_N:K\cap I\neq\emptyset\}
    =
    \binom Nq-\binom{N-q}{q}
    \geq
    \frac{q}{2(q-1)!}N^{q-1}
  \end{equation}
  for all large $N$, because the left-hand side is a polynomial in $N$ with
  leading coefficient $q/(q-1)!$.  Using
  $\sigma_N=((q-1)!)^{1/2}N^{-(q-1)/2}$ gives
  \eqref{eq:single-activity-diverges}.
\end{proof}

The next observation records the other failure mode.  Even after identifying
connected Duhamel cumulants, expanding repeated hyperedges coefficient by
coefficient is too fine a decomposition.

\begin{lemma}[Coefficientwise distinct-support tree bounds are too strong]
  \label{lem:coefficientwise-tree-fails}
  Fix even $q\geq4$.  There is no constant $A_q<\infty$ with the following
  property.  For every family of distinct $q$-sets $E_0,E_1,\ldots,E_s$ and
  every multi-index $a=(a_1,\ldots,a_s)$, the coefficient
  \begin{equation}\label{eq:bad-coefficient-def}
    [z_1^{a_1}\cdots z_s^{a_s}]
    \partial_u\partial_v
    \log\tr\exp\left(
      (u+v)\psi_{E_0}+\sum_{j=1}^s z_j\psi_{E_j}
    \right)\bigg|_{u=v=0}
  \end{equation}
  is bounded in absolute value by
  \begin{equation}\label{eq:false-coefficientwise-bound}
    A_q^{|a|}\,\tau(G),
  \end{equation}
  where $G$ is the intersection graph on
  $\{E_0,E_1,\ldots,E_s\}$ and $\tau(G)$ is its number of spanning trees.
\end{lemma}

\begin{proof}
  Choose a fixed set $S\subset[N]$ with $|S|=q-1$ and distinct points
  $x_0,x_1,\ldots,x_s\notin S$.  Set
  \begin{equation}\label{eq:bad-family-Ej}
    E_j=S\cup\{x_j\},
    \qquad
    A_j=\psi_{E_j}.
  \end{equation}
  Since $|E_i\cap E_j|=q-1$ is odd for $i\neq j$, the operators
  $A_0,A_1,\ldots,A_s$ anticommute pairwise.  Also $A_j^2=\I$.
  Therefore
  \begin{equation}\label{eq:anticommuting-cosh}
    \tr\exp\left(x_0A_0+\sum_{j=1}^s z_jA_j\right)
    =
    \cosh\left(
      \sqrt{x_0^2+\sum_{j=1}^s z_j^2}
    \right).
  \end{equation}
  Writing
  \begin{equation}\label{eq:hq-def}
    h_q(w)=\log\cosh\left(\sqrt w\right),
  \end{equation}
  we get
  \begin{equation}\label{eq:bad-source-derivative-hq}
    \partial_u\partial_v
    \log\tr\exp\left(
      (u+v)A_0+\sum_{j=1}^s z_jA_j
    \right)\bigg|_{u=v=0}
    =
    2h_q'\left(\sum_{j=1}^s z_j^2\right).
  \end{equation}
  The function $h_q'$ has its nearest singularity at a negative real point,
  coming from the first zero of $\cosh$.  Hence, by elementary singularity
  analysis, there are constants $b_q,c_q>0$ such that the coefficient $b_m$ of
  $w^m$ in $h_q'(w)$ satisfies
  \begin{equation}\label{eq:hq-coeff-lower}
    |b_m|\geq b_q c_q^m
  \end{equation}
  for all sufficiently large $m$.
  Take $a_1=\cdots=a_s=2k$ and $m=ks$.  The coefficient of
  $z_1^{2k}\cdots z_s^{2k}$ in \eqref{eq:bad-source-derivative-hq} is
  \begin{equation}\label{eq:multinomial-bad-coeff}
    2b_{ks}\frac{(ks)!}{(k!)^s}.
  \end{equation}
  By Stirling's formula this is at least
  \begin{equation}\label{eq:bad-coeff-growth}
    (c_q')^{ks}s^{ks}
  \end{equation}
  up to subexponential factors in $k$ and $s$.  In the present example the
  intersection graph is the complete graph on $s+1$ vertices, so
  \begin{equation}\label{eq:complete-tree-count}
    \tau(G)=(s+1)^{s-1}.
  \end{equation}
  For any fixed $A_q$, choose $s$ large enough and then $k$ large enough.  The
  lower bound \eqref{eq:bad-coeff-growth} is then larger than
  $A_q^{2ks}(s+1)^{s-1}$, contradicting
  \eqref{eq:false-coefficientwise-bound}.
\end{proof}

\begin{assumption}[Resummed two-replica $L^2$ cluster expansion]
  \label{ass:two-replica-cluster-expansion}
  This is the unproved high-temperature input in the present CLT proof.
  There exist constants $B_q<\infty$ and $\beta_{\rm ce}>0$, depending only on
  $q$, such that the following statements hold for $|\beta|<\beta_{\rm ce}$.

  First, for all $N$ and $I,J\in\calI_N$,
  \begin{equation}\label{eq:cluster-convergence}
    \mathcal C^{[M]}_{I,J}(\beta,g)
    \to
    \mathcal C_{I,J}(\beta,g)
    \qquad
    \text{in }L^2(\dd\gamma(g))
    \text{ as }M\to\infty.
  \end{equation}
  Second, for each $M\geq0$ there is a decomposition in $L^2(\dd\gamma(g))$
  of the finite high-temperature truncation into resummed two-marked cluster
  layers,
  \begin{equation}\label{eq:resummed-layer-decomposition}
    \mathcal C^{[M]}_{I,J}(\beta,g)
    =
    \sum_{p=0}^M
    \mathcal A^{[M]}_{p;I,J}(\beta,g),
  \end{equation}
  where $\mathcal A^{[M]}_{p;I,J}$ is the contribution of Wick-paired connected
  parity clusters with $p$ Gaussian pairs after all repeated-hyperedge
  multiplicities inside each connected cluster have been resummed.  More
  formally, the decomposition is required to be compatible with the exact
  finite Wick formula of Lemma~\ref{lem:finite-two-replica-formula}: the
  $p$-th layer is obtained by collecting all terms with $p$ Wick pairs and by
  summing, before taking absolute values, all diagrams which have the same
  connected support and differ only by repeated occurrences of hyperedges
  inside that support.  Thus $\mathcal A^{[M]}_{0;I,J}$ is the zero-disorder
  covariance layer, supported on the diagonal mark $I=J$, while
  $\mathcal A^{[M]}_{1;I,J}$ is the first annealed correction with one Gaussian
  pair.  These examples are meant only to fix the interpretation of the layers;
  the estimate below is the actual hypothesis.

  The layers satisfy the row-wise $L^2$ estimate
  \begin{equation}\label{eq:resummed-layer-bound}
    \sup_{M\geq0}
    \sup_{N\geq1}
    \sup_{I\in\calI_N}
    \sum_{J\in\calI_N}
    \E\left[
      \left|
        \mathcal A^{[M]}_{p;I,J}(\beta,g)
      \right|^2
    \right]
    \leq
    B_q^p|\beta|^{2p}
    \qquad
    \text{for every }p\geq0.
  \end{equation}
\end{assumption}

\begin{proposition}[Cluster expansion implies two-marked susceptibility]
  \label{prop:cluster-expansion-to-susceptibility}
  Assumption~\ref{ass:two-replica-cluster-expansion} implies
  Assumption~\ref{ass:two-marked-susceptibility}, with any
  \begin{equation}\label{eq:beta-sus-from-ce}
    \beta_{\rm sus}<\min\{\beta_{\rm ce},B_q^{-1/2}\}.
  \end{equation}
  More precisely, if $K\subset(-\beta_{\rm sus},\beta_{\rm sus})$ is compact
  and $\beta_K=\sup_{\beta\in K}|\beta|$, then the susceptibility constant in
  \eqref{eq:two-marked-susceptibility} satisfies
  \begin{equation}\label{eq:susceptibility-constant-from-cluster}
    S_{q,K}
    \leq
    \left(1-B_q^{1/2}\beta_K\right)^{-2}.
  \end{equation}
\end{proposition}

\begin{proof}
  Fix $|\beta|<\min\{\beta_{\rm ce},B_q^{-1/2}\}$.  For the finite truncation,
  Minkowski's inequality
  in the Hilbert space
  $L^2(\dd\gamma(g);\ell^2(\calI_N))$ gives, uniformly in $M,N,I$,
  \begin{align}\label{eq:cluster-to-susceptibility-proof}
    \left(
      \sum_{J\in\calI_N}
      \E\left[
        \left|
          \mathcal C^{[M]}_{I,J}(\beta,g)
        \right|^2
      \right]
    \right)^{1/2}
    &\leq
    \sum_{p=0}^M
    \left(
      \sum_{J\in\calI_N}
      \E\left[
        \left|
          \mathcal A^{[M]}_{p;I,J}(\beta,g)
        \right|^2
      \right]
    \right)^{1/2} \\
    &\leq
    \sum_{p=0}^M
    (B_q^{1/2}|\beta|)^p.
  \end{align}
  The right-hand side is bounded uniformly in $M,N,I$.  Taking $M\to\infty$
  and using \eqref{eq:cluster-convergence}, we obtain
  \begin{equation}\label{eq:susceptibility-from-cluster-limit}
    \sum_{J\in\calI_N}
    \E\left[
      \left|
        \mathcal C_{I,J}(\beta,g)
      \right|^2
    \right]
    \leq
    \left(
      \sum_{p\geq0}(B_q^{1/2}|\beta|)^p
    \right)^2.
  \end{equation}
  Taking the supremum over
  $\beta\in K\subset(-\beta_{\rm sus},\beta_{\rm sus})$, with
  $\beta_{\rm sus}$ satisfying \eqref{eq:beta-sus-from-ce}, gives
  \eqref{eq:susceptibility-constant-from-cluster} and proves
  Assumption~\ref{ass:two-marked-susceptibility}.
\end{proof}

\section{Consistency checks and final conjecture}
\label{sec:checks-conjecture}

\subsection{\texorpdfstring{Small-$\beta$ consistency check}{Small-beta consistency check}}

\begin{lemma}[Trace of a nontrivial Majorana monomial]\label{lem:trace-monomial-zero}
  If $A\subset[N]$ is nonempty, then
  \begin{equation}\label{eq:trace-monomial-zero}
    \tr\left(\prod_{a\in A}\psi_a\right)=0.
  \end{equation}
\end{lemma}

\begin{proof}
  Choose $a\in A$.  By the sign-flip unitary used in
  Lemma~\ref{lem:gauge-invariance}, there is a unitary $U$ such that
  $U\psi_aU^*=-\psi_a$ and $U\psi_bU^*=\psi_b$ for $b\neq a$.  Hence the
  monomial changes sign under conjugation by $U$.  Trace invariance gives the
  claim.
\end{proof}

\begin{lemma}[Trace orthogonality]\label{lem:trace-orth}
  For $I,J\in\calI_N$,
  \begin{equation}\label{eq:trace-orth}
    \tr(\psi_I\psi_J)=\ind_{I=J}.
  \end{equation}
\end{lemma}

\begin{proof}
  If $I=J$, then $\psi_I^2=\I$.  If $I\neq J$, the product
  $\psi_I\psi_J$ is, up to a scalar, a nontrivial Majorana monomial supported
  on $I\triangle J$, so Lemma~\ref{lem:trace-monomial-zero} applies.
\end{proof}

\begin{proposition}[Leading small-$\beta$ fluctuation]\label{prop:small-beta}
  As $\beta\to0$ with $N$ and $g$ fixed,
  \begin{equation}\label{eq:small-beta-expansion}
    L_N(\beta)
    =
    \frac{\beta^2}{2}\tr(H_N^2)
    +
    \mathcal R_N(\beta),
    \qquad
    |\mathcal R_N(\beta)|
    \leq
    8|\beta|^3\norm{H_N}^3,
  \end{equation}
  and
  \begin{equation}\label{eq:tr-H2}
    \tr(H_N^2)=\sigma_N^2\sum_{I\in\calI_N}g_I^2.
  \end{equation}
  Consequently the leading random term in $L_N(\beta)-\E L_N(\beta)$ is
  \begin{equation}\label{eq:leading-random-term}
    \frac{\beta^2}{2}\sigma_N^2
    \sum_{I\in\calI_N}(g_I^2-1),
  \end{equation}
  and
  \begin{equation}\label{eq:small-beta-variance}
    \Var\left[
      N^{q/2-1}
      \frac{\beta^2}{2}\sigma_N^2
      \sum_{I\in\calI_N}(g_I^2-1)
    \right]
    \to
    \frac{((q-1)!)^2}{2q!}\beta^4.
  \end{equation}
\end{proposition}

\begin{proof}
  Let $F(\beta)=\log\tr e^{-\beta H_N}$.  Since $\tr(H_N)=0$ by
  Lemma~\ref{lem:trace-monomial-zero}, we have $F(0)=F'(0)=0$ and
  $F''(0)=\tr(H_N^2)$.  Lemma~\ref{lem:log-partition-derivative-bound}, applied
  with $B_1=B_2=B_3=H_N$, gives
  $|F'''(t)|\leq2^3 3!\norm{H_N}^3$ for all real $t$.  Taylor's formula gives
  the remainder bound in \eqref{eq:small-beta-expansion}.  By
  Lemma~\ref{lem:trace-orth},
  \begin{equation}\label{eq:tr-H2-proof}
    \tr(H_N^2)
    =
    \sigma_N^2\sum_{I,J\in\calI_N}g_Ig_J\tr(\psi_I\psi_J)
    =
    \sigma_N^2\sum_{I\in\calI_N}g_I^2.
  \end{equation}
  This proves \eqref{eq:tr-H2} and \eqref{eq:leading-random-term}.  Since
  $\Var(g_I^2-1)=2$,
  \begin{equation}\label{eq:small-beta-var-proof}
    \Var\left[
      N^{q/2-1}
      \frac{\beta^2}{2}\sigma_N^2
      \sum_I(g_I^2-1)
    \right]
    =
    N^{q-2}\frac{\beta^4}{4}\sigma_N^4 2M_N.
  \end{equation}
  Using $\sigma_N^2=(q-1)!N^{-(q-1)}$ and $M_N\sim N^q/q!$ gives
  \eqref{eq:small-beta-variance}.
\end{proof}

\subsection{Final conjecture}

The following is a self-contained conjectural form of
Assumption~\ref{ass:two-replica-cluster-expansion}.  Fix an even integer $q\geq4$.  For each $N$, let
  $\psi_1,\ldots,\psi_N$ be Majorana generators on the normalized trace space,
  and for $I=\{i_1<\cdots<i_q\}\subset[N]$ put
  $\psi_I=\psi_{i_1}\cdots\psi_{i_q}$ with the standard phase convention for
  which $\psi_I$ is Hermitian and $\psi_I^2=\I$.  Let
  \begin{equation}\label{eq:final-conj-index-set}
    \calI_N=\{I\subset[N]:|I|=q\},
    \qquad
    \sigma_N^2=(q-1)!N^{-(q-1)}.
  \end{equation}
  Let $(g_I)_{I\in\calI_N}$ be independent standard Gaussian variables and set
  \begin{equation}\label{eq:final-conj-hamiltonian}
    H_N(g)=\sigma_N\sum_{I\in\calI_N}g_I\psi_I,
    \qquad
    Z_N(\beta)=\tr e^{-\beta H_N(g)}.
  \end{equation}
  For $I,J\in\calI_N$, define the Kubo connected covariance
  \begin{equation}\label{eq:final-conj-kubo}
    \mathcal C_{I,J}(\beta,g)
    =
    \int_0^1
    \left\{
      \frac{
        \tr\left(\psi_I e^{-s\beta H_N(g)}
        \psi_J e^{-(1-s)\beta H_N(g)}\right)
      }{Z_N(\beta)}
      -
      \frac{\tr(\psi_Ie^{-\beta H_N(g)})}{Z_N(\beta)}
      \frac{\tr(\psi_Je^{-\beta H_N(g)})}{Z_N(\beta)}
    \right\}
    \dd s.
  \end{equation}
  For $m\geq0$ and $K_1,\ldots,K_m\in\calI_N$, define
  \begin{equation}\label{eq:final-conj-kappa}
    \kappa(I,J;K_1,\ldots,K_m)
    =
    \partial_u\partial_v\partial_{t_1}\cdots\partial_{t_m}
    \log\tr\exp\left(
      u\psi_I+v\psi_J+\sum_{a=1}^m t_a\psi_{K_a}
    \right)\bigg|_{u=v=t_1=\cdots=t_m=0}.
  \end{equation}
  For $M\geq0$, put
  \begin{equation}\label{eq:final-conj-truncation}
    \mathcal C^{[M]}_{I,J}(\beta,g)
    =
    \sum_{m=0}^M   \mathcal A_{p;I,J}(\beta,g),
    \end{equation}
    where
    \[
 \mathcal A_{p;I,J}(\beta,g)=     \frac{(-\beta\sigma_N)^m}{m!}
    \sum_{K_1,\ldots,K_m\in\calI_N}
    g_{K_1}\cdots g_{K_m}
    \kappa(I,J;K_1,\ldots,K_m).
  \]

\begin{conjecture}[Self-contained two-replica cluster expansion]
  \label{conj:self-contained-two-replica-cluster}
  There exist constants $B_q<\infty$ and $\beta_{q}>0$, depending only on
  $q$, such that for every $|\beta|<\beta_{q}$ the following hold:
  \begin{equation}\label{eq:final-conj-convergence}
    \lim_{M\to\infty} \sup_{N,I,J\in\calI_N}
    \E\left[
      \left(
        \mathcal C^{[M]}_{I,J}(\beta,g)
        -
        \mathcal C_{I,J}(\beta,g)
      \right)^2
    \right]
    =
    0.
  \end{equation}
  \end{conjecture}
  \begin{conjecture}[Self-contained two-replica cluster expansion]
  \label{conj:self-contained-two-replica-cluster2}
  There exist constants $B_q<\infty$ and $\beta_{q}>0$, depending only on
  $q$, such that for every $|\beta|<\beta_{q}$ the following hold,
  uniformly in $N$: \begin{equation}\label{eq:final-conj-layer-bound}
    \sup_{M\geq0}
    \sup_{N\geq1}
    \sup_{I\in\calI_N}
    \sum_{J\in\calI_N}
    \E\left[
      \left|
        \mathcal A_{p;I,J}(\beta,g)
      \right|^2
    \right]
    \leq
    B_q^p|\beta|^{2p}
    \qquad
    \text{for every }p\geq0.
  \end{equation}
\end{conjecture}

\begin{thebibliography}{GPSZ26}

\bibitem[GPSZ26]{GPSZ2026}
David Gamarnik, Francisco Pernice, Alexander Schmidhuber, and Alexander
Zlokapa.
\newblock The free energy limit of the {SYK} model at high temperature.
\newblock arXiv:2605.02768, 2026.

\bibitem[ACKK25]{ACKK2025}
Eric R. Anschuetz, Chi-Fang Chen, Bobak T. Kiani, and Robbie King.
\newblock Strongly interacting fermions are non-trivial yet non-glassy.
\newblock Physical Review Letters, 135(3):030602, 2025.

\bibitem[Wil67]{Wilcox1967}
R. M. Wilcox.
\newblock Exponential operators and parameter differentiation in quantum physics.
\newblock Journal of Mathematical Physics, 8(4):962--982, 1967.

\end{thebibliography}

\end{document}
